\section{Test}

Questa sezione descrive la strategia di verifica del progetto, con l'obiettivo di mostrare come il sistema sia stato controllato dal punto di vista funzionale, applicativo e, dove possibile, dell'esperienza utente.

\subsection{Obiettivi della fase di test}

La fase di test ha lo scopo di verificare:

\begin{itemize}
	\item correttezza dei flussi applicativi principali;
	\item coerenza delle API esposte dai servizi;
	\item corretto comportamento dell'interfaccia nei casi d'uso principali;
	\item adeguatezza generale dell'esperienza utente.
\end{itemize}

\subsection{Test backend}

In questa sottosezione si possono riportare i test svolti o pianificati sui servizi backend, distinguendo se necessario tra:

\begin{itemize}
	\item test unitari della logica applicativa;
	\item test di integrazione degli endpoint REST;
	\item verifica dei controlli di autorizzazione;
	\item verifica della correttezza dei payload in ingresso e in uscita.
\end{itemize}

Se una parte dei test non e' ancora stata completata, conviene dichiararlo in modo esplicito, indicando quali verifiche risultano gia' disponibili e quali rappresentano lavoro futuro.

\subsection{Test frontend}

Questa sottosezione puo' descrivere le verifiche effettuate sull'interfaccia, ad esempio:

\begin{itemize}
	\item corretto rendering delle viste principali;
	\item aggiornamento dei dati a seguito delle azioni dell'utente;
	\item comportamento delle componenti di configurazione, chatbot e notifiche;
	\item fruibilita' su viewport differenti.
\end{itemize}

\subsection{Test end-to-end dei flussi principali}

E' utile riportare almeno alcuni scenari completi, come:

\begin{itemize}
	\item registrazione o login e accesso alle funzionalita' protette;
	\item creazione di una configurazione e salvataggio;
	\item popolamento del carrello a partire da una configurazione;
	\item aggiornamento del catalogo da parte dell'amministratore e osservazione degli effetti sul sistema.
\end{itemize}

Per ciascuno scenario si possono indicare passi principali, risultato atteso e risultato osservato.

\subsection{Valutazione qualitativa con utenti}

Se sono stati svolti test informali o osservazioni su utenti reali o simulati, questa e' la sede giusta per riportarli. In particolare si possono discutere:

\begin{itemize}
	\item facilità d'uso percepita del configuratore;
	\item comprensibilita' del flusso di acquisto;
	\item utilita' percepita di chatbot e notifiche;
	\item eventuali criticita' emerse durante la prova.
\end{itemize}

\subsection{Limiti dell'attivita' di test}

In questa parte conviene esplicitare i limiti della verifica svolta, ad esempio presenza di test manuali predominanti, copertura non completa o parti ancora non implementate. Questa trasparenza rende il report piu' credibile e aiuta a contestualizzare i risultati.

\subsection{Sintesi dei risultati}

La sezione puo' chiudersi con una sintesi dei principali esiti della verifica, evidenziando quali obiettivi risultano gia' soddisfatti e quali miglioramenti restano pianificati per l'evoluzione del progetto.